%% ============================================================
%% Contexto del Caso — Proyecto Final  [v2, compacta]
%% Cóndor (fintech) — Introducción al Business Analytics — 06278-ECO
%% Universidad Icesi | Periodo 2026-01
%% ============================================================
\documentclass[11pt,letterpaper]{article}

\usepackage[utf8]{inputenc}
\usepackage[T1]{fontenc}
\usepackage[spanish]{babel}
\usepackage[top=2.2cm, bottom=2.0cm, left=2.4cm, right=2.4cm]{geometry}
\usepackage{parskip}
\usepackage{booktabs}
\usepackage{tabularx}
\usepackage{array}
\usepackage{enumitem}
\usepackage{fancyhdr}
\usepackage{titlesec}
\usepackage[table]{xcolor}
\usepackage{tcolorbox}
\usepackage{amssymb}
\usepackage{tikz}
\usepackage{longtable}
\usepackage{hyperref}
\usepackage{lastpage}
\tcbuselibrary{skins,breakable}

%% ─── Paleta del curso ────────────────────────────────────────
\definecolor{azulcurso}{HTML}{1F4E79}
\definecolor{azulmedio}{HTML}{2E75B6}
\definecolor{naranja}{HTML}{E87722}
\definecolor{gristabla}{HTML}{F2F5F9}

\hypersetup{colorlinks=true, linkcolor=azulcurso, urlcolor=azulmedio}
\setlength{\emergencystretch}{3em}

%% ─── Secciones (compactas) ───────────────────────────────────
\titleformat{\section}{\normalsize\bfseries\color{azulcurso}}{\thesection.}{0.5em}{}[\vspace{1pt}{\color{naranja}\hrule height 1pt}\vspace{2pt}]
\titleformat{\subsection}{\small\bfseries\color{azulmedio}}{\thesubsection}{0.5em}{}
\titlespacing{\section}{0pt}{10pt}{4pt}
\titlespacing{\subsection}{0pt}{6pt}{2pt}

%% ─── Encabezado y pie ────────────────────────────────────────
\pagestyle{fancy}
\fancyhf{}
\fancyhead[L]{\footnotesize\itshape\color{azulcurso} Caso Cóndor --- Proyecto Final}
\fancyhead[R]{\footnotesize Business Analytics $|$ Icesi 2026-01 $|$ Pág.\ \thepage/\pageref{LastPage}}
\renewcommand{\headrulewidth}{0.4pt}

%% ─── Columnas y listas ───────────────────────────────────────
\newcolumntype{L}[1]{>{\raggedright\arraybackslash}p{#1}}
\setlength{\tabcolsep}{5pt}
\renewcommand{\arraystretch}{1.15}
\setlist[itemize]{label=\textbullet, leftmargin=1.3em, itemsep=1pt, topsep=2pt}

%% ─── Cajas de resaltado ──────────────────────────────────────
\newtcolorbox{cajaclave}{enhanced, breakable, colback=gristabla, colframe=azulmedio,
  boxrule=0.7pt, arc=2pt, left=7pt, right=7pt, top=4pt, bottom=4pt,
  fonttitle=\bfseries\small\color{azulcurso}, title={Idea clave}}
\newtcolorbox{cajaalerta}{enhanced, breakable, colback=naranja!8, colframe=naranja,
  boxrule=0.7pt, arc=2pt, left=7pt, right=7pt, top=4pt, bottom=4pt,
  fonttitle=\bfseries\small\color{naranja!85!black}, title={Regla de oro: no usar información del futuro}}

\newcommand{\var}[1]{\texttt{\detokenize{#1}}}
\newcommand{\target}{{\color{naranja}$\bigstar$}}

%% ============================================================
\begin{document}
\thispagestyle{fancy}

%% ─── Encabezado del documento (sin portada aparte) ───────────
\begin{center}
  {\color{naranja}\rule{\textwidth}{1.2pt}}\\[4pt]
  {\LARGE\bfseries\color{azulcurso} Caso Cóndor}\\[2pt]
  {\normalsize\color{azulmedio} Contexto del Proyecto Final --- Introducción al Business Analytics (06278-ECO)}\\[3pt]
  {\footnotesize Universidad Icesi\ $|$\ Periodo 2026-01\ $|$\ Fecha de corte de los datos: 31 de marzo de 2026}\\[3pt]
  {\color{naranja}\rule{\textwidth}{1.2pt}}
\end{center}

\vspace{2pt}

%% ============================================================
\section{La empresa y el encargo}

\textbf{Cóndor} es una \textit{fintech} (empresa ficticia) con sede en Colombia y operación incipiente en Perú y Ecuador. Opera una plataforma de \textbf{pagos y crédito}: una billetera (\textit{wallet}) para que las personas paguen con QR y tarjeta, transfieran, recarguen y pidan microcréditos, y un lado de comercios que aceptan esos pagos.

Cinco años de crecimiento acelerado dejaron a Cóndor con \textbf{datos abundantes pero desarticulados}. La Gerencia sospecha que la empresa crece en usuarios pero pierde dinero por los flancos ---fraude, mora de crédito, clientes que se vuelven inactivos, campañas poco efectivas--- y no logra cuantificar dónde ni por qué. \textbf{Ustedes son el equipo de analistas contratado} para elegir uno de esos frentes, formular una pregunta de negocio y responderla con las herramientas del curso.

\begin{cajaclave}
El proyecto \textbf{no} evalúa código perfecto. Evalúa que planteen una buena pregunta, elijan la tarea analítica adecuada (clasificación, regresión o segmentación), la ejecuten con orden e \textbf{interpreten} los resultados para que alguien sin conocimientos técnicos pueda decidir y actuar.
\end{cajaclave}

%% ============================================================
\section{Los frentes de trabajo y sus rutas}

Cada grupo elige \textbf{un} frente. La tabla muestra la pregunta típica de cada área y la ruta analítica natural; varias preguntas admiten más de un enfoque.

\begin{center}
\footnotesize
\begin{tabularx}{\textwidth}{L{3.0cm} X L{2.9cm}}
  \toprule
  \rowcolor{gristabla}
  \textbf{Frente (área)} & \textbf{Pregunta típica} & \textbf{Ruta natural} \\
  \midrule
  Retención (Growth) & ¿Qué clientes se van a volver inactivos? & Clasificación \\
  Finanzas / Planeación & ¿Cuánto va a gastar un cliente el próximo trimestre? & Regresión \\
  Segmentación (Marketing) & ¿Qué grupos de clientes existen y cómo se comportan? & Segmentación \\
  Fraude y Seguridad & ¿Qué transacciones son fraudulentas? & Clasificación \\
  Riesgo de Crédito & ¿Qué solicitudes de crédito caerán en mora? & Clasificación / Regresión \\
  Marketing (campañas) & ¿Quién va a responder a una campaña? & Clasificación \\
  \bottomrule
\end{tabularx}
\end{center}

%% ============================================================
\section{La estructura de los datos}

Cóndor les entrega \textbf{una base principal} con lo esencial y \textbf{tres anexos} con detalle opcional. Todas las tablas se conectan por \var{cliente_id}.

\vspace{2pt}
\begin{center}
\begin{tikzpicture}[
    every node/.style={font=\scriptsize},
    base/.style={rectangle, rounded corners=3pt, draw=azulcurso, fill=azulmedio!12,
                 text width=4.0cm, align=center, minimum height=1.2cm, very thick},
    anexo/.style={rectangle, rounded corners=3pt, draw=azulmedio, fill=gristabla,
                  text width=3.2cm, align=center, minimum height=0.95cm, thick},
    link/.style={draw=naranja, thick}]
  \node[base] (b) at (0,0) {\textbf{base\_clientes}\\ \scriptsize base principal\\ \scriptsize (una fila por cliente)};
  \node[anexo] (t) at (-4.6,-2.5) {\textbf{anexo\_transacciones}\\ \scriptsize (una fila por transacción)};
  \node[anexo] (c) at ( 0,-2.5)   {\textbf{anexo\_creditos}\\ \scriptsize (una fila por solicitud)};
  \node[anexo] (m) at ( 4.6,-2.5) {\textbf{anexo\_campanas}\\ \scriptsize (una fila por envío)};
  \draw[link] (b) -- (t);  \draw[link] (b) -- (c);  \draw[link] (b) -- (m);
  \node at (0,-1.35) {\scriptsize\itshape se unen por cliente\_id};
\end{tikzpicture}
\end{center}

\vspace{-2pt}
\begin{itemize}
  \item \textbf{Ruta fácil.} Trabajen \textbf{solo} con la base principal (\var{base_clientes}). Ya trae todo calculado a nivel cliente y alcanza para las tres rutas.
  \item \textbf{Ruta avanzada.} Sumen uno o más anexos para hacer su propio \textit{feature engineering} o abordar preguntas de fraude, crédito o campañas. Es más exigente y se reconoce en la rúbrica.
\end{itemize}

%% ============================================================
\section{Base principal: \texorpdfstring{\var{base_clientes}}{base\_clientes}}

Una fila por cliente. Reúne demografía, comportamiento (RFM), uso de la app, soporte y crédito, más \textbf{dos variables objetivo} ya listas. Con esta sola tabla se pueden hacer las tres rutas. A continuación se explica cada variable.

\vspace{3pt}
{\footnotesize
\setlength{\LTpre}{2pt}\setlength{\LTpost}{2pt}
\begin{longtable}{L{4.4cm} L{1.6cm} L{9.2cm}}
  \toprule
  \rowcolor{gristabla} \textbf{Variable} & \textbf{Tipo} & \textbf{Descripción} \\
  \midrule
  \endfirsthead
  \multicolumn{3}{r}{\footnotesize\itshape\color{azulcurso} (continúa de la página anterior)}\\[1pt]
  \toprule
  \rowcolor{gristabla} \textbf{Variable} & \textbf{Tipo} & \textbf{Descripción} \\
  \midrule
  \endhead
  \midrule
  \multicolumn{3}{r}{\footnotesize\itshape\color{azulcurso} continúa en la página siguiente…}\\
  \endfoot
  \bottomrule
  \endlastfoot

  \multicolumn{3}{l}{\cellcolor{azulmedio!10}\textbf{\textit{Identificación y demografía}}}\\
  \var{cliente_id}        & id         & Identificador único del cliente. Es la llave para unir la base con los anexos. \\
  \var{edad}              & entero     & Edad del cliente, en años. \\
  \var{genero}            & categórica & Género del cliente: F, M u Otro. \\
  \var{ciudad}            & categórica & Ciudad de residencia declarada. \\
  \var{departamento}      & categórica & Departamento o región de residencia. \\
  \var{ocupacion}         & categórica & Ocupación del cliente: empleado, independiente, estudiante o pensionado. \\
  \var{ingreso_declarado} & continua   & Ingreso mensual que el cliente declaró al registrarse, en pesos colombianos (COP). \\
  \var{kyc_nivel}         & ordinal    & Nivel de verificación de identidad (nivel\_1, nivel\_2 o nivel\_3): a mayor nivel, más datos del cliente fueron validados. \\
  \var{canal_adquisicion} & categórica & Cómo llegó el cliente a Cóndor: orgánico, ads\_meta, ads\_google, referido o influencer. \\
  \var{antiguedad_meses}  & entero     & Meses transcurridos desde que el cliente se registró hasta la fecha de corte. \\
  \addlinespace
  \multicolumn{3}{l}{\cellcolor{azulmedio!10}\textbf{\textit{Comportamiento transaccional (RFM)}}}\\
  \var{recencia_dias}         & entero   & Días desde la última transacción del cliente hasta el corte. Es la \textbf{R} de RFM: menos días significa un cliente más reciente. \\
  \var{frecuencia_tx}         & entero   & Número total de transacciones que hizo el cliente. Es la \textbf{F} de RFM. \\
  \var{monto_total}           & continua & Suma de todos los montos que transó el cliente, en COP. Es la \textbf{M} de RFM. \\
  \var{monto_promedio}        & continua & Monto promedio por transacción, en COP. \\
  \var{ticket_mediano}        & continua & Mediana del monto de sus transacciones, en COP. A diferencia del promedio, no se distorsiona con transacciones muy grandes. \\
  \var{pct_tx_internacional}  & continua & Proporción de sus transacciones que fueron internacionales, entre 0 y 1. \\
  \var{num_categorias_mcc}    & entero   & Cantidad de rubros distintos en los que ha comprado (variedad de consumo). \\
  \var{canal_preferido}       & categórica & Canal que el cliente usa con más frecuencia: app\_qr, tarjeta, pse o link. \\
  \addlinespace
  \multicolumn{3}{l}{\cellcolor{azulmedio!10}\textbf{\textit{Uso de la app, soporte y saldo}}}\\
  \var{num_sesiones_ult30d}       & entero   & Número de sesiones que el cliente abrió en la app en los últimos 30 días antes del corte. \\
  \var{duracion_sesion_promedio}  & continua & Duración promedio de una sesión, en segundos. Queda \textbf{vacía} si el cliente no tuvo sesiones recientes. \\
  \var{os_principal}              & categórica & Sistema operativo principal del cliente: Android o iOS. \\
  \var{num_tickets}               & entero   & Número de tickets de soporte que ha abierto el cliente. \\
  \var{csat_promedio}             & continua & Satisfacción promedio con el soporte, en escala de 1 a 5. Queda \textbf{vacía} si el cliente nunca abrió un ticket. \\
  \var{gasto_promedio_mensual}    & continua & Gasto mensual histórico promedio del cliente, en COP. \\
  \var{saldo_promedio_ult3m}      & continua & Saldo promedio en la cuenta durante los últimos 3 meses, en COP. \\
  \addlinespace
  \multicolumn{3}{l}{\cellcolor{azulmedio!10}\textbf{\textit{Crédito}}}\\
  \var{tiene_credito} & binaria  & 1 si el cliente ha solicitado al menos un crédito; 0 si nunca lo ha hecho. \\
  \var{score_buro}    & continua & Puntaje de buró de crédito (de 300 a 850) de su última solicitud. Queda \textbf{vacía} si el cliente nunca pidió crédito. \\
  \addlinespace
  \multicolumn{3}{l}{\cellcolor{naranja!12}\textbf{\textit{Variables objetivo (lo que un modelo intenta predecir)}}}\\
  \target~\var{abandono} & binaria & \textbf{Ruta de clasificación.} Vale 1 si el cliente se vuelve inactivo (deja de usar Cóndor) y 0 si sigue activo. Está definida para \textbf{todos} los clientes. \\
  \target~\var{gasto_proximo_trim} & continua & \textbf{Ruta de regresión.} Gasto total del cliente, en COP, durante el trimestre siguiente a la fecha de corte. \\
\end{longtable}
\par}

\begin{cajaclave}
Algunas columnas tienen faltantes \textbf{con sentido}: \var{score_buro} está vacía para quien nunca pidió crédito, \var{csat_promedio} para quien no abrió tickets y \var{duracion_sesion_promedio} para quien no tuvo sesiones recientes. Diagnosticar y decidir qué hacer con esos faltantes es parte de la Entrega~2.
\end{cajaclave}

%% ============================================================
\section{Anexos (detalle opcional)}

Cada anexo se une a la base por \var{cliente_id} y habilita preguntas más específicas. Su variable objetivo está marcada con \target.

{\footnotesize
\textbf{\var{anexo_transacciones}} --- una fila por transacción. Permite detección de fraude, construir RFM propio o segmentar con más detalle.\\
\textit{Columnas:} \var{transaccion_id}, \var{cliente_id}, \var{comercio_id}, \var{timestamp}, \var{monto}, \var{tipo} (compra, transferencia, retiro, recarga, pago\_servicio), \var{canal}, \var{rubro}, \var{es_internacional}, \var{estado} (aprobada/rechazada), \target\ \var{etiqueta_fraude}.

\vspace{3pt}
\textbf{\var{anexo_creditos}} --- una fila por solicitud de crédito. Permite predecir la mora o el monto recuperado de la cartera.\\
\textit{Columnas:} \var{credito_id}, \var{cliente_id}, \var{fecha_solicitud}, \var{monto_aprobado}, \var{plazo_meses}, \var{tasa}, \var{score_buro}, \var{destino}, \target\ \var{default_90d} (cayó en mora de 90+ días), \var{dias_mora}, \target\ \var{monto_recuperado}.

\vspace{3pt}
\textbf{\var{anexo_campanas}} --- una fila por envío de campaña a un cliente. Permite predecir la respuesta a campañas de marketing.\\
\textit{Columnas:} \var{campana_id}, \var{cliente_id}, \var{fecha_envio}, \var{canal} (push/email/sms), \var{tipo_oferta} (cashback, descuento, upgrade\_credito, genérico), \var{costo}, \target\ \var{convertido}.
\par}

\begin{cajaalerta}
No usen como predictor ninguna variable que solo se conoce \textbf{después} del evento que quieren predecir. Por ejemplo, \var{dias_mora} y \var{monto_recuperado} \textbf{no} sirven para predecir \var{default_90d}: al momento de aprobar el crédito todavía no existen. Cuando una variable ``predice demasiado bien'', sospechen de esta fuga de información.
\end{cajaalerta}

%% ============================================================
\section{Cómo elegir la ruta y reglas de los datos}

\textbf{Primero la pregunta, luego la técnica.} Decidan qué quieren responder y de ahí se desprende la ruta:
\begin{itemize}
  \item \textbf{Clasificación} (predecir una categoría): \var{abandono} en la base, o \var{etiqueta_fraude}, \var{default_90d}, \var{convertido} en los anexos.
  \item \textbf{Regresión} (predecir un número): \var{gasto_proximo_trim} en la base, o \var{monto_recuperado} en el anexo de créditos.
  \item \textbf{Segmentación} (agrupar sin objetivo): usen el bloque RFM de la base. Un segmento solo sirve si se puede \textbf{describir} y \textbf{accionar} (``clientes jóvenes de alto gasto y baja antigüedad'', no ``grupo 3'').
\end{itemize}

\textbf{Dos reglas.} (1) \textit{Fecha de corte:} los predictores se miden hasta el 31 de marzo de 2026; los objetivos se refieren a lo que ocurre en o después de esa fecha. (2) \textit{Entrenamiento y prueba:} separen los datos, entrenen solo con la parte de entrenamiento y midan el desempeño en la de prueba, comparándolo siempre contra un \textbf{baseline} (la clase mayoritaria en clasificación; la media o mediana en regresión).

\textbf{Cuidado con el grano al unir tablas.} La base tiene una fila por cliente, pero \var{anexo_transacciones} tiene muchas por cliente. Para llevar transacciones al nivel de cliente hay que \textbf{agregar primero} (contar, sumar, promediar) y luego unir; unir sin agregar duplica filas.

%% ============================================================
\section{Cargar los datos en R}

Las cuatro tablas están publicadas como \texttt{.csv} y se cargan con \texttt{read.csv()} apuntando a su URL:

\vspace{2pt}
{\small\ttfamily
\noindent url\_base <- "https://eduard-martinez.github.io/databases/ba/condor/"\\[1pt]
\noindent clientes <- read.csv(paste0(url\_base, "base\_clientes.csv"))\\
\noindent \# anexos (opcionales, para la ruta avanzada):\\
\noindent transacciones <- read.csv(paste0(url\_base, "anexo\_transacciones.csv"))
}
\vspace{3pt}

Para la ruta fácil basta con \var{base_clientes} (es liviana). El \var{anexo_transacciones} es el más grande; si solo necesitan una parte, fíltrenlo apenas lo carguen.

\vspace{4pt}
\begin{center}
{\footnotesize\itshape\color{azulcurso} Cóndor es una empresa ficticia creada con fines pedagógicos. Cualquier parecido con una empresa real es coincidencia.}
\end{center}

\end{document}
